\subsection*{最初に必要な用語}
\addcontentsline{toc}{subsection}{最初に必要な用語}

\par\smallskip\noindent\refstepcounter{table}\label{tab:ch06-01}\textbf{表 \thetable: 第06章 ソフトウェアエンジニアリング運用 表01}\par\smallskip
\begin{longtable}{L{35mm}L{115mm}}
\toprule
用語 & 初学者向けの説明 \\
\midrule
ソフトウェアエンジニアリング運用 & ソフトウェアを配備し、動かし、監視し、障害や変更に対応し、利用者へ安定したサービスを提供し続ける活動である。 \\
運用担当技術者 & 運用サービス、配備、監視、障害対応、容量管理、データ保護などを実行または自動化する技術者である。 \\
開発運用連携（DevOps） & 開発、保守、運用の間の壁を低くし、共通のプロセスとツールで素早く安全にソフトウェアを改善する考え方である。 \\
インフラストラクチャのコード化（IaC） & サーバ、ネットワーク、設定、環境構築などを手作業ではなくコードとして表し、再現可能に管理する方法である。 \\
プラットフォーム工学 & 開発者が自分で環境構築、配備、監視などを使えるように、内部向けの共通基盤を作る活動である。 \\
サイト信頼性エンジニアリング（SRE） & 可用性、性能、遅延、容量、緊急対応などを、ソフトウェア工学の方法で改善する運用の考え方である。 \\
サービス水準合意（SLA） & サービスの可用性、応答時間、復旧時間などについて、提供者と利用者の間で合意した水準である。 \\
インシデント & サービスの正常な運用を妨げる出来事である。障害、性能劣化、誤設定、セキュリティ事象などが含まれる。 \\
問題 & 一つまたは複数のインシデントの背後にある原因、または原因候補である。 \\
テレメトリ & ソフトウェアやインフラから収集する運用データである。ログ、実行トレース、CPU使用量、メモリ使用量、利用者操作などが含まれる。 \\
\bottomrule
\end{longtable}
